% ============================================================
%  Openbook / Formelsammlung Template — strikt nach my_dhbw
%  Stil: 2-spaltig, kompakt, scriptsize Tabellen, dünne Regeln.
%  Renderer: lualatex (zweimal, wg. Unicode-Inhalt)
% ============================================================
\documentclass[10pt, a4paper]{article}

\usepackage[ngerman]{babel}
\usepackage{fontspec}
\setmainfont[
  Path=/usr/share/fonts/TTF/,
  Extension=.ttf,
  BoldFont=DejaVuSans-Bold,
  ItalicFont=DejaVuSans-Oblique,
  BoldItalicFont=DejaVuSans-BoldOblique
]{DejaVuSans}
\setsansfont[
  Path=/usr/share/fonts/TTF/,
  Extension=.ttf,
  BoldFont=DejaVuSans-Bold,
  ItalicFont=DejaVuSans-Oblique,
  BoldItalicFont=DejaVuSans-BoldOblique
]{DejaVuSans}
\setmonofont[Path=/usr/share/fonts/TTF/,Extension=.ttf]{DejaVuSansMono}

\usepackage[top=0.8cm, bottom=0.8cm, left=0.8cm, right=0.8cm, footskip=0.4cm]{geometry}
\usepackage{amsmath, amssymb}
\usepackage{multicol}
\usepackage{xcolor}
\usepackage{titlesec}
\usepackage{enumitem}
\usepackage{fancyhdr}
\usepackage{lastpage}
\usepackage{microtype}
\usepackage{etoolbox}

\usepackage{booktabs}
\usepackage{array}
\usepackage{tabularx}
% ltablex/longtable bewusst NICHT geladen — kollidiert mit multicols.
\usepackage{colortbl}
\usepackage[most]{tcolorbox}
\usepackage{listings}
\usepackage[normalem]{ulem}
\usepackage{adjustbox}
\usepackage[colorlinks=false]{hyperref}

\renewcommand{\familydefault}{\sfdefault}

% ── Theme (my_dhbw accent) ────────────────────────────────────
\definecolor{accent}{RGB}{0, 60, 113}
\definecolor{primaryblue}{RGB}{0, 60, 113}
\definecolor{codebg}{RGB}{245, 245, 245}
\definecolor{figurebg}{RGB}{232, 241, 255}
\definecolor{tablehintbg}{RGB}{240, 255, 240}
\definecolor{solutionbg}{RGB}{242, 255, 242}
\definecolor{rulegray}{RGB}{180, 180, 180}
\definecolor{tableheadbg}{RGB}{0, 60, 113}

% ── Header/Footer — wie my_dhbw formelsammlung.tex ────────────
\pagestyle{fancy}
\fancyhf{}
\renewcommand{\headrulewidth}{0pt}
\renewcommand{\footrulewidth}{0.3pt}
\fancyfoot[L]{\tiny\color{accent}\nichttitle}
\fancyfoot[R]{\tiny\color{accent}\thepage\,/\,\pageref{LastPage}}

% ── Typografie — wie formelsammlung.tex ───────────────────────
\titleformat{\section}
    {\color{accent}\footnotesize\bfseries}{}{0pt}{}
    [\vspace{-2pt}\color{accent!50}\rule{\linewidth}{0.4pt}]
\titleformat{\subsection}
    {\scriptsize\bfseries\color{accent!85!black}}{}{0pt}{}{}
\titleformat{\subsubsection}
    {\scriptsize\itshape\color{accent!60!black}}{}{0pt}{}{}
\titlespacing*{\section}{0pt}{4pt}{1pt}
\titlespacing*{\subsection}{0pt}{3pt}{0pt}
\titlespacing*{\subsubsection}{0pt}{2pt}{0pt}

\setlength{\parindent}{0pt}
\setlength{\parskip}{1pt}
\renewcommand{\arraystretch}{1.0}
\setlist{noitemsep, topsep=0pt, parsep=0pt, leftmargin=1.1em}

\setlength{\columnsep}{6mm}
\setlength{\columnseprule}{0.2pt}

% Globaler scriptsize-Body. Tabellen zusätzlich auf scriptsize zwingen.
\AtBeginEnvironment{tabularx}{\scriptsize}
\AtBeginEnvironment{tabular}{\scriptsize}

% ── Boxen (kompakt) ───────────────────────────────────────────
\newtcolorbox{figurebox}[1]{
  colback=figurebg, colframe=accent!50,
  title={\scriptsize Abbildung: #1}, fonttitle=\scriptsize\bfseries,
  fontupper=\scriptsize, sharp corners,
  boxsep=2pt, left=3pt, right=3pt, top=2pt, bottom=2pt,
  before skip=3pt, after skip=3pt
}
\newtcolorbox{tablehintbox}[1]{
  colback=tablehintbg, colframe=green!50!black!50,
  title={\scriptsize Tabelle: #1}, fonttitle=\scriptsize\bfseries,
  fontupper=\scriptsize, sharp corners,
  boxsep=2pt, left=3pt, right=3pt, top=2pt, bottom=2pt,
  before skip=3pt, after skip=3pt
}
\newtcolorbox{solutionbox}[1]{
  colback=solutionbg, colframe=green!60!black!60,
  title={\scriptsize #1}, fonttitle=\scriptsize\bfseries,
  fontupper=\scriptsize, sharp corners,
  boxsep=2pt, left=3pt, right=3pt, top=2pt, bottom=2pt,
  before skip=3pt, after skip=3pt
}

\lstset{
  basicstyle=\ttfamily\tiny,
  backgroundcolor=\color{codebg},
  breaklines=true,
  frame=none,
  xleftmargin=2pt, xrightmargin=2pt,
  aboveskip=2pt, belowskip=2pt,
  keepspaces=true
}

\providecommand{\nichttitle}{Formelsammlung}
\providecommand{\nichtsubtitle}{}

\renewcommand{\nichttitle}{Relationen, Algebra, Diskrete Mathematik}
\renewcommand{\nichtsubtitle}{Formelsammlung}


\begin{document}

\begin{center}
{\small\bfseries\color{accent}\nichttitle}
\end{center}
\vspace{-6pt}

\scriptsize
\begin{multicols}{2}

% ── BODY ──────────────────────────────────────────────────────

\section{Logik}


\subsection{Grundlagen}
\begin{itemize}
  \item \textbf{Aussage}: eindeutig w/f. Variablen: a,b,c (nicht w/f). \textbf{Belegung}: Variablen → \{w,f\}.
  \item \textbf{Atomar}: nicht weiter zerlegbar. Verknüpft mit ∧, ∨, ¬.
\end{itemize}


\subsection{Operatoren}

{\small
\begin{adjustbox}{max width=\linewidth}
{\scriptsize
\begin{tabular}{lll}
\toprule
\textbf{Sym} & \textbf{Name} & \textbf{Regel} \\
\midrule
∧ & Konjunktion & wahr ⇔ beide wahr \\
∨ & Disjunktion & wahr ⇔ mind. 1 wahr \\
¬ & Negation & – \\
→ & Implikation & = ¬a∨b; falsch nur (w,f) \\
↔ & Äquivalenz & = (a∧b)∨(¬a∧¬b) \\
\bottomrule
\end{tabular}
}
\end{adjustbox}
}


\begin{itemize}
  \item \textbf{Neutrale}: w für ∧ (a∧w=a); f für ∨ (a∨f=a)
  \item \textbf{Vorrang}: ( )<¬<∧<∨<→<↔ (0–5); ∧ wie · weglassen; gleich → links nach rechts
\end{itemize}


\subsection{Syntax / Semantik}
\begin{itemize}
  \item Syntax-Regeln: (1) xᵢ, w, f Formeln; (2) (a∧b),(a∨b) Formeln; (3) ¬a Formel
  \item → und ↔ NICHT in Syntax (durch ∧/∨/¬ darstellbar)
  \item \textbf{Erfüllbar}: ∃ Belegung F=w · \textbf{Tautologie}: alle w · \textbf{Widerspruch}: alle f
  \item F Tautologie ↔ ¬F widerspruchsvoll
\end{itemize}


\subsection{Logische Gesetze}

{\small
\begin{adjustbox}{max width=\linewidth}
{\scriptsize
\begin{tabular}{ll}
\toprule
\textbf{Name} & \textbf{Form} \\
\midrule
Idempotenz & a∧a↔a; a∨a↔a \\
Kommutativ & a∧b↔b∧a \\
Assoziativ & (a∧b)∧c↔a∧(b∧c) \\
Absorption & a∧(a∨b)↔a; a∨(a∧b)↔a \\
Distributiv & a∧(b∨c)↔(a∧b)∨(a∧c); a∨(b∧c)↔(a∨b)∧(a∨c) \\
Doppelte Neg. & ¬(¬a)↔a \\
DeMorgan & ¬(a∧b)↔¬a∨¬b; ¬(a∨b)↔¬a∧¬b \\
Kontraposition & a→b ↔ ¬b→¬a (≠ ¬a→¬b) \\
\bottomrule
\end{tabular}
}
\end{adjustbox}
}



\subsection{Normalformen}
\begin{itemize}
  \item \textbf{Literal}: Variable oder negierte Variable
  \item \textbf{K} (Konj.term): x₁∧…∧xₙ · \textbf{D} (Disj.term): x₁∨…∨xₙ
  \item \textbf{DN}: K₁∨…∨Kₘ · \textbf{KN}: D₁∧…∧Dₘ
  \item \textbf{kDN/kKN}: jeder Term enthält ALLE Variablen → Min/Maxterme; eindeutig
\end{itemize}


\subsubsection{Verfahren DN/KN}
\begin{enumerate}
  \item ↔, → ersetzen (¬,∨,∧)
  \item DeMorgan bis ¬ nur vor Variablen
  \item Distributiv (DN: ∧ über ∨; KN: ∨ über ∧)
  \item Idempotenz; Neutrale: ¬a∧a=f, ¬a∨a=w, a∧w=a, a∨f=a
\end{enumerate}


\subsubsection{Verfahren kDN/kKN}
\begin{itemize}
  \item Fehlende Var a: K → K∧(a∨¬a); D → D∨(a∧¬a); dann Distributiv + Idempotenz
  \item Bsp: a → a∧(b∨¬b) → ab∨a¬b
  \item \textbf{Sortierung}: lexikographisch in K/D; ¬x:=0, x:=1 → Binärzahl aufsteigend
\end{itemize}


\subsection{Logisches Schließen / Resolution}
\begin{itemize}
  \item \{(a→b),(b→c)\} ⊨ (a→c) — Voraussetzungen ⟹ Schlussfolgerung
  \item \textbf{Beweis durch Widerspruch}: zeige (A₁∧…∧Aₙ∧¬b) widerspruchsvoll. Voraussetzung: KN
  \item \textbf{Klausel}: Disjunktion von Literalen als Menge
  \item \textbf{Resolvente}: A∋a, B∋¬a ⟹ res\{A,B\}=(A∖\{a\})∪(B∖\{¬a\}); pro Schritt 1 Literal
  \item Komplementäre Literale in Resolvente → verwerfen
  \item \textbf{□} (leere Klausel) = Widerspruch bewiesen
  \item Einzelne Klausel nicht resolvierbar; Klauseln bleiben in Menge
  \item Kein □, kein Schritt mehr → Behauptung falsch
\end{itemize}

\begin{quote}
\textbf{Merke:} Resolution für ↔-Vermutungen unhandlich → SAT.
\end{quote}


\subsection{Prädikatenlogik}
\begin{itemize}
  \item \textbf{P(x₁,…,xₙ)} n-stelliges Prädikat · ∀ Allquantor · ∃ Existenzquantor · Funktionsterme liefern Individuen
  \item "Alle Menschen sterblich": ∀x(Mensch(x)→sterblich(x))
  \item "kein Mensch springt >20m": ¬∃x(Mensch(x)∧springt\_weiter(x,20))
\end{itemize}


\section{Mengen}


\subsection{Begriff}
\begin{itemize}
  \item Cantor: Zusammenfassung unterscheidbarer Objekte. Großbuchstaben für Mengen, klein für Elemente
  \item \textbf{Aufzählend} \{1,2,3\} · \textbf{Beschreibend} \{x|A(x)\} · \textbf{Venn}
  \item Ungültig: \{1,2,3,3\}; andere Klammern
  \item a∈M, a∉M; |M| Mächtigkeit
  \item |∅|=0; |\{∅\}|=1; |P(∅)|=1; |P(\{∅\})|=2
\end{itemize}


\subsection{Leere Menge ∅}
\begin{itemize}
  \item = \{\}; Teilmenge JEDER Menge; nicht Element jeder Menge (nur jeder Potenzmenge)
\end{itemize}


\subsection{Gleichheit / Teilmenge}
\begin{itemize}
  \item Extensionalität: M=N ↔ ∀x(x∈M↔x∈N)
  \item T⊂M echt · T⊆M unecht · M⊇T Obermenge
\end{itemize}


\subsection{Zahlenmengen}

{\small
\begin{adjustbox}{max width=\linewidth}
{\scriptsize
\begin{tabular}{ll}
\toprule
\textbf{Menge} & \textbf{Inhalt} \\
\midrule
ℕ & \{1,2,3,…\} \\
ℕ₀ & \{0,1,2,…\} \\
ℤ & \{…,−1,0,1,…\} \\
ℚ & \{p/q \textbackslash{} \\
ℝ & Dezimalbrüche a₀.a₁a₂…, a₀∈ℤ, aᵢ∈\{0,…,9\} \\
ℂ & komplex \\
\bottomrule
\end{tabular}
}
\end{adjustbox}
}


\begin{itemize}
  \item [a,b] abgeschl.; (a,b) offen; (a,b], [a,b) halboffen
\end{itemize}


\subsection{Potenzmenge / Komplement}
\begin{itemize}
  \item \textbf{P(M)}=\{x|x⊆M\}; |P(M)|=2\textasciicircum{}|M|
  \item \textbf{C\_E(M)}=M̄=E∖M (M⊆E)
\end{itemize}


\subsection{Mengenalgebra}

{\small
\begin{adjustbox}{max width=\linewidth}
{\scriptsize
\begin{tabular}{ll}
\toprule
\textbf{Op} & \textbf{Definition} \\
\midrule
∩ & \{x \textbackslash{} \\
∪ & \{x \textbackslash{} \\
∖ & \{x \textbackslash{} \\
Disjunkt & A∩B=∅ \\
× & \{(x,y) \textbackslash{} \\
\bottomrule
\end{tabular}
}
\end{adjustbox}
}


\begin{itemize}
  \item Gesetze analog Logik (Idempotenz, Komm., Assoz., Absorption, Distributiv, Doppeltes Komplement, DeMorgan)
\end{itemize}


\section{Algebraische Strukturen}

\begin{itemize}
  \item \textbf{Algebr. Struktur}: Menge + Verknüpfungen + Axiome; v: M×M→M
\end{itemize}


\subsection{Gruppe (G,+)}
\begin{enumerate}
  \item Abgeschlossen: a+b∈G
  \item Assoziativ
  \item Neutrales e: a+e=a
  \item Inverses a\textit{: a+a}=e
\end{enumerate}

\begin{itemize}
  \item \textbf{Kommutativ}: zusätzlich a+b=b+a
  \item (ℝ,+),(ℤ,+) komm. Gruppen · (ℕ₀,+) keine Gruppe (kein Inverses) · (ℝ∖\{0\},∗) komm. Gruppe
\end{itemize}


\subsection{Körper (K,+,∗)}
\begin{enumerate}
  \item + und ∗ definiert
  \item (K,+) komm. Gruppe
  \item (K∖\{0\},∗) komm. Gruppe
  \item Distributiv: a∗(b+c)=(a∗b)+(a∗c)
  \item (ℝ,+,∗) Körper
\end{enumerate}


\section{Relationen}


\subsection{Begriff}
\begin{itemize}
  \item 2-stellig: R⊆M₁×M₂ · R⊆M² wenn M=M₁=M₂
  \item (a,b)∈R, aRb, R(a,b)
  \item <,>,≤,≥,=,≠,≈ sind Relationen
  \item n-stellig: R⊆M₁×…×Mₙ; DB-Relation = n-stellig, Tabellenkopf = Schema
\end{itemize}


\subsection{Eigenschaften}

{\small
\begin{adjustbox}{max width=\linewidth}
{\scriptsize
\begin{tabular}{ll}
\toprule
\textbf{Eigensch.} & \textbf{Bedingung} \\
\midrule
Reflexiv & ∀x: xRx \\
Irreflexiv & ∄x: xRx \\
Symmetrisch & xRy → yRx \\
Asymmetrisch & xRy → ¬yRx \\
Antisymmetrisch & xRy ∧ yRx → x=y \\
Transitiv & xRy ∧ yRz → xRz \\
Linear & ∀x,y: xRy ∨ yRx \\
Konnex & x≠y → xRy ∨ yRx \\
\bottomrule
\end{tabular}
}
\end{adjustbox}
}


\begin{quote}
\textbf{Merke:} asymmetrisch ⇒ antisymmetrisch.
\end{quote}


\subsection{Äquivalenz / Ordnung}
\begin{itemize}
  \item \textbf{Äquivalenz}: reflexiv + symmetrisch + transitiv → Äquivalenzklassen
\end{itemize}


{\small
\begin{adjustbox}{max width=\linewidth}
{\scriptsize
\begin{tabular}{lll}
\toprule
\textbf{Ordnung} & \textbf{Irreflexiv} & \textbf{Reflexiv} \\
\midrule
Quasi & irreflex.+trans. & reflex.+trans. \\
Halb & + asymmetr. & + antisymm. (≤) \\
Voll & + konnex (<) & + linear (≤) \\
\bottomrule
\end{tabular}
}
\end{adjustbox}
}


\begin{quote}
\textbf{Merke:} "≥" ist KEINE reflex. Halbordnung (nicht antisymm.).
\end{quote}


\section{Abbildungen / Funktionen}

\begin{itemize}
  \item f: D→W, x↦f(x); spez. 2-stellige Relation
  \item \textbf{Linksvollständig}: jedes x∈D zugeordnet · \textbf{Rechtseindeutig}: max. 1 Bild
\end{itemize}


{\small
\begin{adjustbox}{max width=\linewidth}
{\scriptsize
\begin{tabular}{lll}
\toprule
\textbf{Eigensch.} & \textbf{Definition} & \textbf{Pfeile pro y} \\
\midrule
Injektiv & f(x₁)=f(x₂)→x₁=x₂ & max. 1 \\
Surjektiv & ∀y∈W ∃x: f(x)=y & ≥1 \\
Bijektiv & inj.∧surj. & genau 1; invertierbar \\
\bottomrule
\end{tabular}
}
\end{adjustbox}
}


\begin{itemize}
  \item f:ℝ→ℝ, x² → weder · f:ℝ⁺→ℝ⁺, x² → bijektiv (D, W entscheiden!)
\end{itemize}


\section{Relationale Datenbanken}

\begin{itemize}
  \item Relation=Tabelle · Tupel=Datensatz · Schema=Kopf
  \item \textbf{Primärschlüssel}: eindeutig + minimal
  \item \textbf{Fremdschlüssel}: in anderer Relation PK
  \item Schlechter Entwurf → Redundanz/Anomalien; Lösung: Normalisierung + FK
\end{itemize}


\subsection{Relationenalgebra}

{\small
\begin{adjustbox}{max width=\linewidth}
{\scriptsize
\begin{tabular}{ll}
\toprule
\textbf{Op} & \textbf{Bedeutung} \\
\midrule
∩,∪,∖ & nur kompatibles Schema; ∪ ohne Duplikate \\
× & ohne Kompat., jedes mit jedem \\
σ\_F(r) & Selektion (Zeilen) — SQL \texttt{where} \\
Π\_Y(r) & Projektion (Spalten Y), keine Duplikate — SQL \texttt{select distinct} \\
⋈ & Join über gleiche Spalten via PK/FK \\
Natural Join & Spalten gleich Name+Inhalt \\
\bottomrule
\end{tabular}
}
\end{adjustbox}
}


\begin{itemize}
  \item SQL inner join: \texttt{from A,B where A.k=B.k}; weitere: left, right, outer, equi, semi
\end{itemize}


\section{Boolesche Algebra}

\textbf{(M; ⊗, ⊕, ∼)}: |M|≥2; ⊗,⊕: M×M→M; ∼: M→M



\subsection{Axiome}

{\small
\begin{adjustbox}{max width=\linewidth}
{\scriptsize
\begin{tabular}{lll}
\toprule
\textbf{Gesetz} & \textbf{Form 1} & \textbf{Form 2} \\
\midrule
Komm. & a⊗b=b⊗a & a⊕b=b⊕a \\
Distrib. & a⊗(b⊕c)=(a⊗b)⊕(a⊗c) & a⊕(b⊗c)=(a⊕b)⊗(a⊕c) \\
Neutrale & ∃1: a⊗1=a & ∃0: a⊕0=a \\
Komplement & a⊗∼a=0 & a⊕∼a=1 \\
\bottomrule
\end{tabular}
}
\end{adjustbox}
}


\begin{quote}
\textbf{Dualität}: ⊗↔⊕, 0↔1; Beweise gelten dual.
\end{quote}


\subsection{Abgeleitet}

{\small
\begin{adjustbox}{max width=\linewidth}
{\scriptsize
\begin{tabular}{lll}
\toprule
\textbf{Name} & \textbf{Form 1} & \textbf{Form 2} \\
\midrule
Idempotenz & a⊗a=a & a⊕a=a \\
Assoz. & (a⊗b)⊗c=a⊗(b⊗c) & (a⊕b)⊕c=a⊕(b⊕c) \\
Absorp. & a⊗(a⊕b)=a & a⊕(a⊗b)=a \\
Doppelt & ∼(∼a)=a & — \\
DeMorgan & ∼(a⊗b)=∼a⊕∼b & ∼(a⊕b)=∼a⊗∼b \\
Neutr.komp. & ∼0=1 & ∼1=0 \\
0/1 & a⊗0=0 & a⊕1=1 \\
\bottomrule
\end{tabular}
}
\end{adjustbox}
}



\subsubsection{Beweis a⊗0=0}
\begin{enumerate}
  \item a⊗∼a=0
  \item a⊗(∼a⊕0)=0
  \item (a⊗∼a)⊕(a⊗0)=0
  \item 0⊕(a⊗0)=0
  \item a⊗0=0
\end{enumerate}


\subsection{Modelle}
\begin{itemize}
  \item \textbf{Mengenalgebra}: P(M); ⊗≙∩, ⊕≙∪, ∼≙C; Neutrale M, ∅
  \item \textbf{Schaltalgebra}: B=\{0,1\}; offen=0/zu=1
\end{itemize}


{\small
\begin{adjustbox}{max width=\linewidth}
{\scriptsize
\begin{tabular}{lll}
\toprule
\textbf{Op} & \textbf{Schaltung} & \textbf{Lampe wenn} \\
\midrule
a∧b & Reihe & beide zu \\
a∨b & Parallel & mind. 1 zu \\
a′ & Ruhekontakt & offen \\
\bottomrule
\end{tabular}
}
\end{adjustbox}
}



{\small
\begin{adjustbox}{max width=\linewidth}
{\scriptsize
\begin{tabular}{llll}
\toprule
\textbf{a} & \textbf{b} & \textbf{a∧b} & \textbf{a∨b} \\
\midrule
0 & 0 & 0 & 0 \\
0 & 1 & 0 & 1 \\
1 & 0 & 0 & 1 \\
1 & 1 & 1 & 1 \\
\bottomrule
\end{tabular}
}
\end{adjustbox}
}



{\small
\begin{adjustbox}{max width=\linewidth}
{\scriptsize
\begin{tabular}{ll}
\toprule
\textbf{a} & \textbf{a′} \\
\midrule
0 & 1 \\
1 & 0 \\
\bottomrule
\end{tabular}
}
\end{adjustbox}
}



\subsection{Boolesche Terme / Funktion}
\begin{itemize}
  \item Syntax: (1) Konst./Var.; (2) A,B Terme → A′,(A∧B),(A∨B); (3) endlich
  \item Vorrang: ( ), ′, ∧, ∨ (1–3); ∧ weglassbar
  \item f: Bⁿ→B; 2ⁿ Belegungen; 2\textasciicircum{}(2ⁿ) Funktionen bei n Vars
\end{itemize}


\subsubsection{Aus Wertetabelle}
\begin{itemize}
  \item \textbf{kDN}: Zeilen mit f=1; Minterm (a=0→a′, a=1→a); Minterme ∨; z.B. 010 → a′bc′
  \item \textbf{kKN}: Zeilen mit f=0; Maxterm (a=0→a, a=1→a′); ∧-verbinden; z.B. 010 → a∨b′∨c
\end{itemize}


\subsection{KV-Diagramm (Karnaugh-Veitch)}
\begin{itemize}
  \item Achsen Gray-Code (00,01,11,10) — 1 Bit Änderung
  \item Nur Einsen eintragen; max. Einserblöcke (auch über Kanten)
  \item Variable konstant in Block → bleibt; wechselt → eliminiert
  \item \textbf{d} (don't-care): frei nutzbar
  \item Bsp 4-Var Block über Kante b′d′: f=b′d′; c=0,d=1 + a=1,d=1: f=ad∨c′d
\end{itemize}

\begin{quote}
\textbf{Merke:} Skaliert nicht (2³²≈4,3·10⁹). Alternative: Quine-McCluskey.
\end{quote}


\section{Schaltnetze}

\textbf{F: Bⁿ → Bᵐ}



\subsection{Halbaddierer}

{\small
\begin{adjustbox}{max width=\linewidth}
{\scriptsize
\begin{tabular}{llll}
\toprule
\textbf{a} & \textbf{b} & \textbf{s} & \textbf{ü} \\
\midrule
0 & 0 & 0 & 0 \\
0 & 1 & 1 & 0 \\
1 & 0 & 1 & 0 \\
1 & 1 & 0 & 1 \\
\bottomrule
\end{tabular}
}
\end{adjustbox}
}


\begin{itemize}
  \item s = a′b ∨ ab′ = a⊕b
  \item ü = ab
\end{itemize}


\subsection{Volladdierer}

{\small
\begin{adjustbox}{max width=\linewidth}
{\scriptsize
\begin{tabular}{lllll}
\toprule
\textbf{a} & \textbf{b} & \textbf{ü₁} & \textbf{s} & \textbf{ü₂} \\
\midrule
0 & 0 & 0 & 0 & 0 \\
0 & 0 & 1 & 1 & 0 \\
0 & 1 & 0 & 1 & 0 \\
0 & 1 & 1 & 0 & 1 \\
1 & 0 & 0 & 1 & 0 \\
1 & 0 & 1 & 0 & 1 \\
1 & 1 & 0 & 0 & 1 \\
1 & 1 & 1 & 1 & 1 \\
\bottomrule
\end{tabular}
}
\end{adjustbox}
}


\begin{itemize}
  \item s = a⊕b⊕ü₁
  \item ü₂ = ab ∨ aü₁ ∨ bü₁
\end{itemize}


\subsection{Spülmaschine}
\begin{itemize}
  \item f(a,b,c)=1 wenn Tür offen ODER (Wasser zu niedrig UND Heizung an)
\end{itemize}


\section{Graphen}


\subsection{Grundlagen}
\begin{itemize}
  \item \textbf{G=(V,E)} bzw. \textbf{G=(V,E,W)}: Knoten, Kanten, Gewichte
  \item Gewicht: Kosten/Distanz/Wahrscheinlichkeit; gerichtet → richtungsabhängig
  \item \textbf{Knotengrad}: \# angrenzende Kanten; Position irrelevant
\end{itemize}


\subsection{Typen}

{\small
\begin{adjustbox}{max width=\linewidth}
{\scriptsize
\begin{tabular}{ll}
\toprule
\textbf{Typ} & \textbf{Merkmal} \\
\midrule
(Un)gerichtet & Kantenrichtung \\
Gewichtet & Werte \\
Multigraph & Mehrfachkanten \\
Leer & keine Kanten \\
Vollständig & alle verbunden \\
Baum & zyklenfrei \\
Kreis & \textbackslash{} \\
Stern & 1 Zentrum \\
Bipartit & 2 Mengen, Kanten nur dazwischen \\
\bottomrule
\end{tabular}
}
\end{adjustbox}
}



\subsection{Darstellung}
\begin{itemize}
  \item \textbf{Adjazenzmatrix} A (n×n): aᵢⱼ≠0 wenn Kante; ungerichtet → symmetr.; Multigraph ohne Gewicht: aᵢⱼ=\#Kanten
  \item \textbf{Adjazenzliste}: pro Knoten Nachbarn; gerichtet 2 Listen (Vor-/Nachfolger); effizient bei spärlichen
\end{itemize}


\subsection{Pfade / Zyklen}
\begin{itemize}
  \item \textbf{Pfad}: Kantensequenz, jeder Knoten max. 1×
  \item \textbf{Zyklus}: geschlossen, Kanten unterschiedlich; Bäume zyklenfrei
  \item \textbf{Eulersch}: jede Kante genau 1×
  \item \textbf{Hamiltonsch}: jeder Knoten genau 1×
\end{itemize}


\section{Kürzeste Pfade}

\begin{itemize}
  \item Pfad mit minimaler Gewichtssumme
  \item \textbf{Brute Force}: alle Touren, optimal aber exponentiell (Bsp 7K/12E → 29 Pfade)
  \item \textbf{Greedy/Best-first}: kleinste Kante; suboptimal, Sackgassen
\end{itemize}


\subsection{Dijkstra}
\begin{itemize}
  \item s → alle anderen
  \item Variablen: d, π, α/previous, u, w(u,v)
  \item KEINE negativen Gewichte; ineffizient für Punkt-zu-Punkt (Abbruch ab visited)
\end{itemize}


\subsubsection{Init}
\begin{lstlisting}
∀v: v.d=∞, v.π=unvisited, v.α=none
s.d=0
\end{lstlisting}



\subsubsection{Update}
\begin{lstlisting}
update(u,v,w):
  if v.d > u.d+w(u,v):
    v.d=u.d+w(u,v); v.α=u
\end{lstlisting}



\subsubsection{Hauptschleife}
\begin{lstlisting}
u=s
while ∃ unvisited (oder Ziel):
  u.π=visited
  for unvisited Nachbar v: update(u,v,w)
  u = unvisited mit min d
\end{lstlisting}



\subsection{A*}
\begin{itemize}
  \item s→z mit Heuristik
  \item Priorität \textbf{p = d + h(v,z)}; Knotenwahl nach min p (nicht d)
  \item Falsche/überschätzende h → suboptimal
\end{itemize}


\subsubsection{Init}
\begin{lstlisting}
∀v: v.d=∞, v.p=∞, v.π=unvisited, v.α=none
s.d=0; s.p=h(s,z)
\end{lstlisting}



\subsubsection{Update}
\begin{lstlisting}
v.d > u.d+w → v.d=…; v.α=u; v.p=v.d+h(v,z)
\end{lstlisting}



\subsubsection{Hauptschleife}
\begin{lstlisting}
while z.π==unvisited:
  u.π=visited
  for unvisited Nachbarn v: update(...)
  u = unvisited mit min p
\end{lstlisting}



\subsection{Vergleich}

{\small
\begin{adjustbox}{max width=\linewidth}
{\scriptsize
\begin{tabular}{ll}
\toprule
\textbf{Algo} & \textbf{Eigenschaft} \\
\midrule
Dijkstra & 1 Start; ger.+unger.; keine neg. Gewichte \\
Bellman–Ford & 1 Start; neg. Gewichte erlaubt \\
A* & Start+Ziel; Heuristik \\
Floyd–Warshall & alle Paare; gerichtet, dicht \\
Johnson & alle Paare; gerichtet, spärlich \\
\bottomrule
\end{tabular}
}
\end{adjustbox}
}



\subsection{R-Implementation}
\begin{itemize}
  \item \texttt{graph[x,y]} = Gewicht/0
  \item \texttt{distances}(Inf), \texttt{visited}(FALSE), \texttt{previous}(NA), \texttt{priorities}(Inf, A*), \texttt{heuristic}
\end{itemize}


\subsubsection{\texttt{dijkstra(graph,start,end=FALSE)}}
\begin{lstlisting}
distances=rep(Inf,n); visited=rep(FALSE,n); previous=rep(NA,n)
distances[start]=0; previous[start]=0; current=start
while(any(!visited)):
  visited[current]=TRUE
  if(current==end) break
  for i Nachbarn:
    if graph[current,i]!=0 && !visited[i] && distances[i]>distances[current]+graph[current,i]:
      distances[i]=distances[current]+graph[current,i]; previous[i]=current
  current = unvisited mit min distances; sonst -1
\end{lstlisting}

\begin{itemize}
  \item \texttt{end=FALSE} → Pfade zu allen
\end{itemize}


\subsubsection{A*}
\begin{itemize}
  \item + \texttt{priorities[start]=heuristic[start]}
  \item Update setzt zusätzlich \texttt{priorities[i]=distances[i]+heuristic[i]}
  \item Knotenwahl: min \texttt{priorities[i]} über \texttt{!visited[i]}
\end{itemize}


\subsubsection{\texttt{getpath(previous,start,end)}}
\begin{lstlisting}
i=end; path=end
while(i!=0):
  ii=previous[i]; i=ii
  if(i!=0) path=c(ii,path)
return path
\end{lstlisting}

\begin{itemize}
  \item Termination: \texttt{previous[start]=0}
\end{itemize}


\subsubsection{\texttt{allpaths(graph)}}
\begin{lstlisting}
for s in 1:n: results[[s]]=dijkstra(graph,start=s,end=FALSE)
\end{lstlisting}

\begin{itemize}
  \item n Aufrufe (statt n²)
\end{itemize}


\subsubsection{Ravensburg–Immenstaad (6K, ungerichtet)}

{\footnotesize
\begin{adjustbox}{max width=\linewidth}
{\scriptsize
\begin{tabular}{lllllll}
\toprule
\textbf{} & \textbf{1} & \textbf{2} & \textbf{3} & \textbf{4} & \textbf{5} & \textbf{6} \\
\midrule
1 & 0 & 1 & 0 & 0 & 4 & 0 \\
2 & 1 & 0 & 1 & 0 & 1 & 0 \\
3 & 0 & 1 & 0 & 1 & 0 & 0 \\
4 & 0 & 0 & 1 & 0 & 6 & 1 \\
5 & 4 & 1 & 0 & 6 & 0 & 0 \\
6 & 0 & 0 & 0 & 1 & 0 & 0 \\
\bottomrule
\end{tabular}
}
\end{adjustbox}
}


\begin{itemize}
  \item start=1, end=6; A*-Heuristik c(1,0,0,0,1000,0) → Knoten 5 gemieden
\end{itemize}


\subsubsection{Visualisierung (R)}
\begin{itemize}
  \item \texttt{network}, \texttt{ggnet}, \texttt{GGally}
  \item \texttt{network(m,directed=F,names.eval="weights",ignore.eval=FALSE)}
  \item \texttt{ggnet2(...,mode="fruchtermanreingold",size=10,label=T,edge.label="weights",arrow.gap=0.05)}
  \item Pfadknoten rot; \texttt{set.seed(43)}
\end{itemize}


\section{TSP}

\begin{itemize}
  \item n Städte alle genau 1×, zurück zum Start, min. Tourlänge
  \item Graph: V=Städte, E=Strecken, W=Längen
\end{itemize}


\subsection{Anwendungen}
\begin{itemize}
  \item Logistik (VRP), Mikrochip-Design/Leiterplatten, Genom-Sequenzierung (Knoten=DNA-Stränge)
\end{itemize}


\subsection{Verwandt}
\begin{itemize}
  \item TSP ohne Rückkehr · VRP (mehrere) · Traveling Thief (TSP+Knapsack) · stochastisch/dynamisch
\end{itemize}


\subsection{Lösungsverfahren}
\begin{itemize}
  \item \textbf{Brute Force}: alle Touren; optimal, exponentiell
  \item \textbf{NNA (Nearest Neighbour)}: Greedy ab Startknoten
\begin{enumerate}
  \item Start
  \item Nächste = günstigste Kante zu unbesuchtem
  \item Bis alle besucht
  \item Startabhängig, suboptimal; Verbesserung: jeden Start probieren (×n, weiter ohne Garantie)
\end{enumerate}
  \item \textbf{2-Opt}: Lokale Suche zur Tour-Verbesserung
\begin{enumerate}
  \item Approx. Startlösung (NNA)
  \item Iterativ Teilstück invertieren falls kürzer → übernehmen
  \item Nächster Tausch
  \item Approximativ
\end{enumerate}
  \item Weitere (nicht behandelt): lineares Programmieren, evolutionäre Algorithmen
\end{itemize}






% ──────────────────────────────────────────────────────────────

\end{multicols}
\end{document}
